\section{Related work}
\label{sec:related-work}

\subsection{Shipboard registration and antenna pointing}

Shipboard marine radar already uses acquisition-time navigation data when motion within a rotation matters. McCann and Bell calibrated azimuth, range, and timing offsets by maximizing the sharpness of shipborne X-band images accumulated using timestamped radar and navigation data \cite{McCann2018OffsetCalibration}. Lund et al. interpolated position and heading to individual radar pulses for sea-ice drift retrieval and treated heave, pitch, and roll as smaller than the 7.5 m range resolution for their application \cite{Lund2018SeaIceDrift}. Hessner et al. likewise georeferenced pulses with time-matched navigation data and relied on the broad elevation beam of the radar aboard \emph{Polarstern} to accommodate its typical pitch and roll \cite{Hessner2019Reliability}. On vessels under 15 m, Simpson et al. aligned radar-derived wave fields using wave and buoy signatures and identified antenna-height variation as a possible source of amplitude error \cite{Simpson2020Forecasting}. Together, these studies establish pulse-time registration and show that the importance of vertical motion depends on the measurement scale and platform.

Burnett et al. measured the effect of compensating motion distortion on spinning-radar navigation \cite{Burnett2021MotionDistortion}. Chen et al. showed how time-varying attitude affects the observability of position and attitude biases for mobile radar \cite{Chen2014Registration}. Xie et al. studied vertical motion of shipborne X-band radar in significant-wave-height estimation \cite{Xie2019WaveCompensation}. These studies evaluate navigation, registration, or wave retrieval. Here the design quantity is the change in the number of attempted buoy visits that yield an accurate position when pointing, timestamp, or both are changed.

In coastal radar odometry, Jang et al. improved estimates of vessel rotation by matching successive radar images with maritime scene features \cite{Jang2024MaritimeOdometry}. That correction acts on recorded data. By comparison, Mar et al. evaluated ship-motion compensation that changes the direction of a phased-array beam toward an airborne target \cite{Mar2013Motion}. A rotating radar observing a buoy likewise couples registration and illumination to the same vessel motion, but the two interventions enter at different stages. Registration changes the position assigned to an acquired echo. Pointing changes the signal available to the detector during the visit.

\subsection{Crest clearance and low-elevation target returns}

Studies of marine radar imaging show why the direct ray alone cannot predict received power. Wijaya and van Groesen used geometric shadowing to generate synthetic sea-surface radar images for wave-height retrieval \cite{Wijaya2016Shadowing}. In measurements of deep-water ocean backscatter, Plant and Farquharson found that geometric shadowing poorly described low-grazing-angle dropouts \cite{Plant2012WaveShadowing}. Janssen et al. calculated X-band diffraction around smooth wave crests and found finite illumination within nominal ray shadows and reduced illumination immediately outside them \cite{Janssen2026Diffraction}. Their calculation treats incident illumination of the water surface. A buoy echo also depends on scattering from the target and propagation back to the receiver.

Herselman et al. measured calibrated X-band returns from small craft at medium-to-low grazing angles. Small-boat echoes faded as the craft entered wave troughs, and the return from a boat heading into the waves showed periodicity at the mean wave period \cite{Herselman2008SmallBoat}. Those measurements connect wave phase to the strength of a low-elevation target echo. They also identify multipath and clutter as contributors to variability. Clutter-adaptive radar detectors estimate the local background to set a threshold \cite{Rohling1983RadarCFAR}. Thus a prescribed clear-path margin can isolate the effects of obstruction and pointing, while estimates of detection probability for a particular radar and buoy require measured target and clutter returns.

\subsection{Wave shape and vessel response}

The crest and vessel attitude share the same wave forcing, but their relative phase depends on the hull response. Mounet et al. fitted wave-to-motion transfer functions to inertial data from an AutoNaut unmanned surface vessel. The estimated heave and pitch spectra agreed more closely with measurements than the roll spectra did, even after calibration \cite{Mounet2024Hydrodynamic}. Crest shape introduces a separate effect. Second-order wave interactions sharpen crests and change their height distribution \cite{Tayfun1980Nonlinear,Forristall2000Crest}. McAllister and van den Bremer measured distinct second-order motions for a surface-following buoy and a fixed wave gauge in the same steep wave groups \cite{McAllister2019Lagrangian}. These results motivate treating hull response, intervening crest shape, and buoy motion as separate inputs to path clearance.

The present comparison links these inputs over the same scheduled visits. Processing the same echo with two poses measures the mapping penalty due to the scan timestamp. Recomputing acquisition with level pointing reveals which visits cross the echo threshold when the antenna changes attitude. Counting consecutive failures tests the influence of scan timing on the gaps between usable positions.
